\scanpage{67}
% letterhead
\noindent U N I V E R S I T E\hfill Montpellier, le 28 Juin 1978\\
DES SCIENCES ET TECHNIQUES\\
DU LANGUEDOC\\
\texttt{-{}-{}-{}-{}-{}-{}-{}-{}-{}-{}-{}-}\\
Service de la Scolarité\\
MR / GJ\\
Diffusion\quad : Liste A\\
N° : 4066

% printed form
\uline{AVIS DE SOUTENANCE D'UNE THESE DE DOCTORAT DE 3ème CYCLE}

\uline{( CHIMIE ORGANIQUE PHYSIQUE )}

\uline{Nom et prénom du candidat} :\qquad AURELLY Michèle

\uline{Titre de la thèse} : Application de la corrélation de Hammett à l'étude de l'énolisation des cétones en milieux basique et superbasique.

\uline{Date, heure et lieu de soutenance prévus} :\\
JEUDI 13 JUILLET 1978 à 16 Heures\\
Salle de cours n° 9 - U. S. T. L.

\uline{Nom du Directeur de thèse} :\\
Monsieur le Professeur LAMATY.

\uline{Désignation du laboratoire où la thèse est préparée} :\\
Laboratoire de Chimie Organique Physique.

\uline{Analyse de l'ouvrage} :

Ce travail consiste à cerner de plus près les mécanismes de l'énolisation des cétones par l'étude d'acétophénones et de méthyl-2 phényl-1 butanones-1 substituées, dans des milieux basiques. Les constantes p de Hammett ont été mesurées en fonction de la nature du milieu.

\uline{Durée approximative des recherches ayant abouti à la thèse : 3 ans.}

\uline{Composition du jury proposé} :\\
MM. :\\
LAMATY, Président\\
SCHUE\\
WYLDE\\
MION

Le Secrétaire Général,\\
\handwritten{\unsure{E Siau}}\\
\uline{E. SIAU}
% archive sheet number: 66
